\item \textbf{Problema de repaso.}
Un bloque con una bocina atornillada a él se conecta a un resorte que tiene una constante de resorte $k$ y oscila, como se muestra en la figura P16.41. La masa total del bloque y la bocina es $m$, y la amplitud de movimiento de esta unidad es $A$. La bocina emite ondas sonoras de frecuencia $f$. Determine:
\begin{enumerate}
    \item la frecuencia más alta escuchada por la persona a la derecha de la bocina,
    \item la frecuencia más baja escuchada por la persona a la derecha de la bocina.
    \item Si el nivel sonoro máximo que escucha la persona es $\beta$ cuando está más cerca de la bocina a una distancia $d$, ¿cuál es el nivel sonoro mínimo que escucha el observador?
\end{enumerate}